\documentclass[11pt]{article}
\usepackage[margin=1in]{geometry}
\usepackage{amsmath,amssymb}
\usepackage{fontspec}
\usepackage{polyglossia}
\setmainlanguage{english}
\setotherlanguage{arabic}
\newfontfamily\arabicfont[Script=Arabic]{Amiri}
\newcommand{\ar}[1]{\textarabic{#1}}

\title{Depth Before Derivation: Meta-Level Modification in Orthographic and Morphological Stabilization}
\author{Flyxion}
\date{August 2026}

\begin{document}
\maketitle

\begin{abstract}
\emph{Embodied Constraint and Minimal Assembly} models the stabilization of the
Arabic script as a cost--distinctness optimization: graphemes persist when a
proposed modification's informational gain exceeds its production cost. The
model is powerful but flat. It treats every stabilizing event in the script's
history as an instance of the same operation: a single-level trade between
$C$ and $D$ applied directly to a grapheme. This essay argues that the
script's own history contains at least two operations of different kind, not
merely different degree. Dotting a letter and regularizing a script's
joining rules are not the same type of event. The first modifies a grapheme.
The second modifies the rule that generates a grapheme's possible forms. Read
through an admissibility-dynamics framework, in which a modificatory
operation is classified by its meta-depth $k$, the distinction becomes
visible as a structural fact about the script rather than a stylistic
observation about it. The same distinction reappears in Arabic derivational
morphology, where the traditional measures built on the paradigm root
fa`ala show the pattern with unusual clarity: applying a known measure to a
root is depth-0, while the historical fixing of the measure set itself is
depth-1. This reframing does not replace the cost--distinctness model; it
adds an axis the model does not have, and the resulting picture predicts
which parts of a script's, or a morphology's, history should be
irreversible in a different sense than others.
\end{abstract}

\section{A Flat Model of a Layered History}

\emph{Embodied Constraint} treats the stabilization of a writing system as an
equilibrium condition. A grapheme $g$ persists in a stable configuration when,
for every available augmentation $a$, the ratio of informational gain to
production cost fails to justify adopting $a$:
\[
S = \alpha D - \beta C,
\]
with the script settling near local maxima of $S$. Section 4 of that essay
applies this model to dotting: a minimal graphic addition to a skeletal base
that yields a categorical phonemic distinction, at very low additional cost.
Section 14 applies structurally the same model to a different phenomenon:
the historical replacement of unbounded, context-sensitive ligature forms
with a fixed four-way positional system (isolated, initial, medial, final).
Both are described as instances of the same underlying dynamic, compression
under constraint, cost against distinctness, equilibrium found and held.

The description is not wrong. But it collapses two operations that differ in
what they act \emph{on}. Dotting changes a grapheme. Positional
regularization changes the rule that assigns a grapheme its permissible
shapes. The cost--distinctness formalism, as stated, has no way to mark this
difference; $C(a)$ and $D(g_1,g_2)$ are defined at the level of individual
graphemes throughout, and the model offers no separate variable for
modifications that operate one level up, on the generative rule itself.

\section{Two Kinds of Modification}

Consider the two cases side by side.

\textbf{Dotting.} Before systematic dotting, d\=al and dh\=al, \ar{د} and
\ar{ذ} (to use a rough
transliteration) shared an undotted skeletal base $b$. Dotting introduces an
augmentation $a$ such that $g_1 = b$, $g_2 = a(b)$. The operation acts
directly on the space of admissible graphemes: it adds one new admissible
form. Everything else about the writing system, its joining rules, its
positional behavior, its inventory of augmentation types, is untouched. The
operation reads the current state of the grapheme inventory and writes one
new member into it.

\textbf{Positional regularization.} Before four-position stabilization, the
shape a letter took depended on its specific left and right neighbors,
$F(g_i \mid g_{i-1}, g_{i+1})$, a function whose domain grows combinatorially
in the number of graphemes. Regularization does not add a new grapheme, and
it does not disambiguate two previously conflated forms. It replaces the
function $F$ itself with a lower-complexity function $F(g_i \mid p)$, where
$p$ ranges over four positions rather than over all neighbor pairs. The
object being modified is not a grapheme. It is the rule that had been
generating grapheme forms.

These are not two points on the same axis of \textquotedblleft how much
change occurred.\textquotedblright{} They are modifications with different
targets. One writes into the space of admissible forms under a fixed
generative rule. The other rewrites the generative rule.

\section{A Depth Axis for Orthographic Modification}

This distinction has a name in admissibility-dynamics work developed
elsewhere in this program: meta-depth. A modificatory operation is depth-0,
written $M^{(0)}: R \to R$, if it acts directly on the admissibility
structure $R$, the set of currently reachable or permitted configurations.
An operation is depth-$k$ for $k \geq 1$, written $M^{(k)}: M^{(k-1)} \to
M^{(k-1)}$, if it acts on the machinery that itself modifies $R$, rather than
on $R$ directly. The depth axis is inductive: there is no need for an
$R^{(-1)}$ object, only a base case and a recursive step upward.

Under this axis, dotting is depth-0. It is a write operation on the
admissible-grapheme set $R_t$ at a given moment: one new distinguishable
form is added, everything about how future forms will be generated is left
alone. Positional regularization is depth-1. It is a write operation not on
$R_t$ but on the rule $F$ that determines what $R_{t+1}$ can look like given
$R_t$ and a neighboring context; it changes the generative machinery, not a
product of that machinery.

This is not a difference of magnitude. A very large dotting campaign, applied
to every skeletal ambiguity in the alphabet at once, would still be depth-0:
many individual writes to $R_t$, executed together, but each one a
disambiguation of an existing grapheme pair under an unchanged rule for how
graphemes get their shapes in context. Positional regularization, even
applied to a single letter family first and generalized later, is depth-1
from the moment it changes how shape is determined at all, because after
that point every subsequent grapheme in the affected family is generated by
a different function than before.

\section{The Measures of $f$-`-$l$: Templates as Bounded Generators}

The same depth distinction reappears, more sharply, in Arabic morphology,
where it can be shown rather than only argued for, because the tradition
itself already names its generative apparatus with a fixed exemplar root.
Arabic grammar traditionally illustrates its derivational templates, the
measures or aw\=zan, using the root $f$-`-$l$ (fa\,`a\,la, \ar{فَعَلَ}), whose three
consonants stand in generically for $C_1$, $C_2$, $C_3$ in any root. The
measures are conventionally cited by name using this root directly: Form I
is called \emph{fa`ala}, Form II \emph{fa``ala}, and so on, the pattern name
and the paradigm word being the same string. This is already a small piece
of evidence for the argument to follow: the tradition treats the measures as
a closed, nameable inventory precisely because they function as a fixed
generative rule, not as an open field of possible modification.

Applying fa`ala's own consonants to the ten canonical measures gives:

\begin{center}
\begin{tabular}{llll}
Form & Pattern & Arabic & Gloss (root sense: to do/act) \\ \hline
I & fa`ala & \ar{فَعَلَ} & he did \\
II & fa``ala & \ar{فَعَّلَ} & he made [someone] do; he did intensively \\
III & f\=a`ala & \ar{فَاعَلَ} & he did [something] with/against [someone] \\
IV & 'af`ala & \ar{أَفْعَلَ} & he caused to do \\
V & tafa``ala & \ar{تَفَعَّلَ} & he had himself made to do; did gradually \\
VI & taf\=a`ala & \ar{تَفَاعَلَ} & he pretended to do; they did to one another \\
VII & infa`ala & \ar{اِنْفَعَلَ} & it got done (passive/inchoative) \\
VIII & ifta`ala & \ar{اِفْتَعَلَ} & he did for himself \\
IX & if`alla & \ar{اِفْعَلَّ} & he became [a quality, typically color/defect] \\
X & istaf`ala & \ar{اِسْتَفْعَلَ} & he sought to have [it] done \\
\end{tabular}
\end{center}

Every one of these forms is generated from the same three consonants,
$(C_1,C_2,C_3) = (f,`,l)$, by a different composite operator drawn from a
small set of primitive moves: vowel insertion, gemination of $C_2$
(Form~II), vowel lengthening after $C_1$ (Form~III), prefixation of hamza
(Form~IV) or \emph{ta-} (Forms~V, VI) or \emph{ist-} (Form~X), infixation of
\emph{-t-} after $C_1$ (Form~VIII), or a nasal prefix with final gemination
(Form~VII, IX). In the notation of Appendix D of \emph{Embodied Constraint},
each measure is a composite $T = A_k \circ \dots \circ A_2 \circ A_1$ built
from these primitives, and applying a known measure to a root is a
single write: it inserts one fully specified word $W$ into the space of
forms derivable from $(f,`,l)$, using a generative rule, the finite
measure set $\mathcal{T} = \{T_{\text{I}}, \dots, T_{\text{X}}\}$, that the
operation itself leaves untouched. This is the morphological analogue of
dotting. It is depth-0: a write to $R_t$, the space of admissible derived
words for a given root, under a fixed rule for how derivations happen.

The measure set itself is the depth-1 object. Ten (extended, in some
grammars, to fifteen for rarer patterns) canonical templates is a closed
inventory imposed on what could, in principle, be an unbounded combinatorial
space of primitive-operator compositions: any sequence of vowel insertions,
geminations, prefixations, and infixations applied to a triliteral root
defines some derivable word, but only a small, named, pedagogically fixed
subset of these compositions counts as a productive measure that a speaker
is expected to recognize and extend to new roots. The historical
stabilization of that inventory, whatever process fixed it at approximately
ten forms rather than leaving derivation as an open field of ad hoc
consonant-vowel recombination, is structurally the same event as positional
regularization in the script: a reduction of an unbounded generative space,
$F(g_i \mid g_{i-1}, g_{i+1})$ in one case, arbitrary compositions of
morphological primitives in the other, to a small, closed, teachable rule.
Learning a new root and applying Form II to it is depth-0. Recognizing or
proposing an eleventh productive measure, one that would apply generally
across roots rather than existing as an isolated lexical fossil, would be
depth-1, and this is in fact why such proposals are rare and contested in
the grammatical tradition: a depth-1 change is not merely a bigger version
of a depth-0 one, it is a change to what counts as a legitimate way of
generating a word at all.

This also sharpens a point \emph{Embodied Constraint} makes about
predictive generativity (Section 11.4) without quite stating it in these
terms. The essay observes that speakers can infer the sense of an unattested
derived form by applying a known measure to an unfamiliar root. That
predictive capacity is only available because measure application is
depth-0: it presumes a fixed, already-learned generative rule and simply
executes it on new input. If measure formation were itself a live, per-word
option, generalization from known cases to new ones would not be reliable
in the way the essay describes, since there would be no stable rule to
apply. The productivity of Arabic derivational morphology is, in this
sense, a direct consequence of keeping template formation at depth-1, rare
and historically settled, while leaving template application at depth-0,
frequent and freely generalizable.

\section{Gemination as a Fixed Mark: Shaddah and a Comparative Aside}

Shaddah (\ar{شَدَّة}), the diacritic marking consonant doubling, is depth-0 by the same
test applied to dotting and sukūn: it writes a single explicit fact,
\textquotedblleft this consonant is geminate here,\textquotedblright{} into
$R_t$, the space of admissible syllabic readings, without touching any rule
for how future marks of this kind get assigned. It is, in the vocabulary of
Appendix D, the operator $G(C) = CC$: identity duplication rather than the
introduction of a new graphic primitive, cheaper in $C(a)$ than dotting for
exactly that reason, since no new shape has to be learned, only a
repetition of one already known.

A comparative case outside Arabic sharpens what depth-0 status does and
does not guarantee about a mark's later history. Spanish \~{n} descends from
Latin geminate \emph{nn}, as in \emph{annu} $\to$ \emph{a\~{n}o}. Medieval
scribes abbreviated a doubled letter by writing one instance of it with a
small superscript wave standing in for the second, a purely economical
compression, the same motive Appendix B attributes to ligature simplification
in Arabic: fewer strokes for the same string. At that stage the tilde was
depth-0 in the weakest sense, a shorthand for an already-present doubling,
carrying no phonemic information the doubled letters themselves did not
already carry. Only later, once the geminate nasal had undergone
palatalization into a distinct phoneme, the palatal nasal, did the same mark come
to function the way a dot on d\=al/dh\=al functions: a categorical phonemic
distinguisher, separating \emph{a\~{n}o} from a hypothetical \emph{ano}
rather than merely abbreviating a spelling. The graphic object did not
change. Its function migrated from compression to distinction without any
change in depth, both stages being depth-0 writes, which shows that depth
and function-type, compression versus differentiation, are independent
axes. A mark can move from one function to the other while never leaving
depth-0, so long as neither stage rewrites the rule generating the space of
marks available to writers.

\section{Person, Tense, and Mood: Composing Two Fixed Paradigms}

The measures discussed above generate a derived stem from a root; they do
not by themselves produce a finished, usable verb. A further layer,
inflection for person, gender, number, tense-aspect, and mood, applies on
top of the stem through a second, independently fixed paradigm of prefixes
and suffixes. Holding the root and measure constant at Form I of
$f$-`-$l$, the imperfect indicative alone yields:

\begin{center}
\begin{tabular}{llll}
Form & Arabic & Person/gender & Gloss \\ \hline
'af`alu & \ar{أَفْعَلُ} & 1st singular & I do \\
naf`alu & \ar{نَفْعَلُ} & 1st plural & we do \\
taf`alu & \ar{تَفْعَلُ} & 2nd masc.\ singular / 3rd fem.\ singular & you do / she does \\
yaf`alu & \ar{يَفْعَلُ} & 3rd masc.\ singular & he does \\
yaf`al\=una & \ar{يَفْعَلُونَ} & 3rd masc.\ plural & they do \\
\end{tabular}
\end{center}

Each of these is generated by a fixed paradigm of person-marking affixes,
$\mathcal{P}$, applied to the stem produced by the measure paradigm
$\mathcal{T}$. Neither paradigm modifies the other. Selecting a measure and
selecting a conjugational affix are independent depth-0 writes, each against
its own closed rule set, and a fluent verb form is simply the composition of
both: stem $=$ $T(\text{root})$, word $=$ $P(\text{stem})$. This composability
is itself informative. It is only because $\mathcal{T}$ and $\mathcal{P}$ are
each separately fixed at depth-1 that their composition at depth-0 remains
predictable, a speaker who knows both paradigms can generate or parse a form
never previously encountered.

Mood marking on the imperfect verb makes the connection back to sukūn
explicit rather than analogical. The three moods of the Arabic imperfect are
distinguished only by the final short vowel or its absence: indicative
\emph{yaf`alu} (damma), subjunctive \emph{yaf`ala} (fatḥa), and jussive
\emph{yaf`al} (\ar{يَفْعَلْ}, no vowel at all). The jussive is formed by exactly the
operation Lesson 10 of the Qaedah names Jazm: a sukūn placed on the final
consonant, closing the word where an open syllable had stood. The
traditional Arabic grammatical term for the jussive mood, \emph{majz\=um}
(\ar{مَجْزُوم}),
is built on the same root as \emph{jazm} (\ar{جَزْم}) itself, the term is not merely
analogous to sukūn, it is named for it. The same depth-0 operator that
closes a mid-word syllable in ordinary orthography, marking absence of a
following vowel within a fixed system that otherwise presumes vocalic
continuation, reappears one grammatical level up as the mechanism that
distinguishes an entire mood. Sukūn's status as a null-state write, argued
above to be a boundary case between marking presence and marking absence,
turns out to do exactly the same job whether it closes a syllable inside a
word or closes the final syllable of an inflected verb to signal jussive
mood. The mark did not change. Only the grammatical weight riding on it did.

\section{Waqf: Sukūn at the Boundary of the Utterance}

The pattern identified for the jussive mood, sukūn's closure operator
reused one level up from the syllable to close an entire grammatical
category, recurs a second time at a still larger scope. Arabic nouns in
continuous recitation or connected reading carry case-and-definiteness
endings: an indefinite noun in the nominative takes tanwīn, the doubled
ḍamma pronounced as \emph{-un}, so that the bare consonantal name of the
letter alif, cited indefinitely, is read \emph{alifun} (\ar{أَلِفٌ}), \textquotedblleft an
alif,\textquotedblright{} against the definite \emph{al-alifu} (\ar{الْأَلِفُ}),
\textquotedblleft the alif,\textquotedblright{} where the definite article
replaces the indefinite marking rather than combining with it. This ending
is not optional ornamentation; the presence or absence of tanwīn is what
signals indefiniteness in the first place, doing morphologically what the
English article \emph{a} does syntactically.

At a pause, whether the pause is dictated by the sense of a sentence,
by Qur'ānic waqf convention, or simply by a reader stopping for breath,
this ending is dropped. \emph{Alifun} at pause is read \emph{alif}, its
final vowel and nunation removed regardless of whether the underlying word
was indefinite, definite, nominative, or genitive; the rule governing this
is stated in the Qaedah itself in exactly sukūn's own terms, that the last
letter of a paused word should be made \emph{sakin}. The operation is not a
new mechanism invented for pause. It is the identical closure operator
examined in Section~10.2 of \emph{Embodied Constraint}, applied at the
largest scope encountered so far: not a syllable within a word, not a
mood-marking morpheme on a verb, but the boundary of an entire phrase or
utterance. A reader moving through unvocalized or lightly vocalized
continuous text, a newspaper column being the practical case, uses this
drop precisely as a period or comma is used in scripts that mark
punctuation graphically: the disappearance of the expected final sound is
itself the signal that a unit has closed, doing prosodically what a mark on
the page does visually in other traditions.

Three scopes, one operator. Word-internal sukūn closes a syllable inside a
still-continuing word. Jussive sukūn closes a verb's mood category. Waqf
closes a phrase. In every case the operation is the same depth-0 write
against the same background assumption, that vocalic continuation is the
default and its absence must be explicitly marked, and in every case the
mark's scope, not its identity, is what changes. This is the clearest
demonstration in the corpus examined here that depth and scope are
themselves independent: an operation can remain at depth-0, a direct write
to $R_t$ with no modification to the rule generating $R_t$, while the size
of the unit it closes grows from a single consonant to an entire spoken
sentence.

\section{Hamza and the Selection of a Seat}

Hamza, the glottal stop, is written on one of several carriers depending on
its surrounding vowel environment: on alif, on wāw, on yā', or unsupported
on the line, chosen by a rule that weighs the strength of the vowel before
and after the hamza against a fixed ranking, kasra strongest, ḍamma next,
fatḥa weakest, sukūn weakest of all, and selects the carrier matching the
stronger of the two. Writing \emph{su'āl} with hamza on wāw, or
\emph{ri'āsa} with hamza on yā', is an application of this ranking to a
specific word, a depth-0 write choosing one seat from among the admissible
options given the surrounding vowels, executed against a fixed comparative
rule the operation itself does not touch. \emph{Embodied Constraint}
notes, in Section~10.4, that hamza functions as a boundary marker between
silence and articulated sound. That observation concerns hamza's phonetic
role. The seat-selection rule is a separate, orthographic fact layered on
top of it, and it is this second layer that shows the depth structure:
early Qur'ānic manuscripts are known to have treated hamza inconsistently,
at times omitting it from the rasm entirely, before grammatical
codification fixed both its notation and its seat-selection ranking into
the closed algorithm later scribes and readers inherited as given. The
codification of the ranking was depth-1. Every subsequent application of
it, one hamza, one word, one carrier chosen, is depth-0.

\section{Noon Sākinah and Tanwīn: A Second Partition of the Consonants}

Section~5 of \emph{Embodied Constraint} reads the sun-letter and moon-letter
division as a phonologically motivated partition of the consonant
inventory, coronal versus non-coronal, governing whether the lateral of the
definite article assimilates. A structurally similar but independently
motivated partition governs what happens when noon sākinah or tanwīn is
followed by a consonant. Four treatments are recognized. Idhār, clear
pronunciation with no modification, applies before the throat letters,
the same huroof-e-halqi that block ghunnah in Lesson~5 of the Qaedah.
Idghām, assimilation of the noon into the following consonant, applies
before a small class dominated by nasals and liquids, l\=am, r\=a',
n\=un, m\=im, w\=aw, and y\=a' (\ar{ل ر ن م و ي}), with
ghunnah retained for four of these and dropped for two. Iqlāb, conversion
of the noon sound itself into a meem before articulation, applies solely
before b\=a' (\ar{ب}). Ikhfā, a concealed, partially nasalized pronunciation, applies
before the remaining consonants, the majority of the inventory. This
four-way classification is a second closed taxonomy imposed on the same
twenty-nine consonants the sun-letter division already partitions, keyed
this time not to coronality alone but to a mix of articulatory distance
from noon's own place and manner, nasals and liquids cluster toward
assimilation, the labial b\=a' (\ar{ب}) alone earns full conversion rather than mere
assimilation, sharing noon's nasal manner once converted to meem, and the
throat letters, most distant from noon's place of articulation, block
modification entirely. As with the sun letters, applying the correct
treatment to a specific instance of noon sākinah meeting a specific
following consonant is depth-0, a lookup against a fixed four-way rule.
The historical settling of which consonants belong to which of the four
classes, a classification a beginning reader must simply memorize rather
than derive, is the depth-1 fact making the depth-0 lookups possible at
all. That two independently motivated closed partitions of the same
twenty-nine-letter inventory, one for the definite article, one for noon
sākinah, coexist without needing to agree with one another is itself
evidence that depth-1 stabilization in this system happens locally, rule
by rule, rather than as a single global classification of the consonants
that every subsequent process must respect.

\section{Broken Plurals: A Case the Depth Axis Does Not Close Cleanly}

Not every derivational domain in Arabic resolves as cleanly into a
depth-0/depth-1 pair as the verb measures do. Plural formation offers a
useful complication. Some nouns pluralize by suffixation alone, a sound
plural formed with \emph{-ūn}/\emph{-īn} for masculine or \emph{-āt} for
feminine, a straightforward depth-0 affixation against one of two fixed
suffix paradigms, structurally identical to the person-marking affixes
discussed above. The majority of nouns, however, pluralize by internal
vowel-pattern change, a broken plural, \emph{jam` taks\=ir}: \emph{kit\=ab}
becomes \emph{kutub}, \emph{qalam} becomes \emph{aql\=am},
\emph{rajul} becomes \emph{rij\=al}. This looks, at first, like the nominal
counterpart of the verbal measures examined above, a fixed set of
templates applied to a triliteral base. But the broken-plural template set
is neither small nor fully productive in the way the ten verbal measures
are. Traditional grammars enumerate roughly thirty common patterns, and
critically, which pattern a given singular noun takes is only partially
predictable from its phonological shape; a substantial number of broken
plurals must be learned per noun rather than derived from a rule a speaker
could apply to an unfamiliar root with confidence. This resists a clean
depth assignment. It is not depth-0 in the sense measure-application to
fa`ala is depth-0, since depth-0 there presumes a determinate, executable
rule, and no such rule reliably determines which of thirty patterns a new
noun will take. Nor is it depth-1, since no individual instance of forming
a broken plural modifies the pattern inventory itself. The most accurate
description borrows the interrogative/modificatory distinction developed
elsewhere in this program's admissibility-dynamics work: retrieving the
correct broken-plural pattern for a specific noun functions less like the
execution of a generative rule and more like an interrogative lookup
against a large, only loosely rule-governed stored table, a read operation
over accumulated lexical fact rather than a write executed against a
closed formal grammar. The verb measures earn their productivity, in
Section~11.4 of \emph{Embodied Constraint} and the discussion above, from
sitting cleanly at depth-0 under a closed depth-1 rule. Broken plurals show
what a derivational domain looks like when that condition only partially
holds, and the honest response is not to force a depth number onto it but
to note that not all of Arabic morphology is measure-like, some of it is
closer to memorized association than computed derivation, however much a
given broken plural pattern resembles a measure once it is already known.

\section{What the Cost--Distinctness Model Cannot See}

The consequence of collapsing these into one axis is that the
cost--distinctness model has no principled account of why some orthographic
changes are effectively irreversible in a stronger sense than others. Under
$S = \alpha D - \beta C$, an unwound dotting decision (returning to an
undifferentiated skeletal form) is symmetric with adopting one: both are
single moves along the same $D$-$C$ tradeoff, evaluated at the grapheme
level, and either could in principle be favored again if the weighting
parameters shifted. But undoing positional regularization is not symmetric
with adopting it in the same way. Reverting from four-position allographs
back to fully context-sensitive ligature forms would not just cost more; it
would require re-deriving, or re-inventing, the entire combinatorial
mapping $F(g_i \mid g_{i-1}, g_{i+1})$ that regularization eliminated. The
cost of that reversal is not a point on the same $C$ scale as the cost of
undotting a letter, because it is a cost incurred at the level of the rule
generating the grapheme space, not at the level of a grapheme.

Put in terms of the essay's own resource-cost language: a depth-0 reversal
pays $C(a)$ for a single grapheme. A depth-1 reversal pays for reconstructing
an entire generative function, which is not $C(a)$ summed over graphemes but
a qualitatively different kind of cost, closer to what the admissibility
framework treats as the cost of restoring lost machinery rather than the
cost of restoring a lost state. \emph{Embodied Constraint}'s own appendices
gesture toward this without naming it: Appendix B's treatment of ligature
suppression as \textquotedblleft entropy reduction\textquotedblright{}
correctly identifies that something more than grapheme-level compression is
happening, that the domain of the joining function itself is shrinking from
$\mathcal{O}(n^2)$ to $\mathcal{O}(4n)$, but the essay does not connect this
observation back to its own $S = \alpha D - \beta C$ equilibrium model,
because that model has no slot for a change in the generating rule as
distinct from a change in the generated output.

\section{Sukūn as a Boundary Case}

One diacritic in the same corpus resists clean classification under either
label, and the resistance is informative. Sukūn does not add a new
admissible grapheme, the way dotting does, and it does not rewrite a
generative rule, the way positional regularization does. It marks the
absence of a following vowel: a null state made explicit within a system
that otherwise treats vowel continuation as a default. In the vocabulary
introduced above, sukūn is a depth-0 operation, since it acts directly on
$R_t$, the space of admissible syllabic readings for a given consonant, but
its content is unusual for a depth-0 write: rather than adding a
distinguishable positive form, it declares that no further form follows.
This is closer to what admissibility-dynamics work elsewhere calls a
constitutive operation whose payload is a closure condition rather than a
new member. It is worth flagging as a case that a two-term cost--distinctness
model handles by accident (its $C$ is low, its $D$ from an unmarked
alternative is high, so it fits the same inequality as dotting) without the
model registering that it is doing something categorically different from
adding a new letter.

\section{Conclusion}

\emph{Embodied Constraint}'s central claim, that scripts stabilize under
recurring pressure toward informational distinctness at bounded cost, does
not need to be abandoned to notice that the pressure operates on at least
two different objects across the script's documented history: the grapheme
inventory itself, and the rule that generates grapheme forms from context.
Treating both under a single $S = \alpha D - \beta C$ equilibrium obscures a
distinction that the script's own evolution makes visible when the objects
of modification are tracked rather than assumed uniform. Dotting operates
within a fixed rule for what a letter can look like. Positional
regularization operates on that rule. The difference is not one of scale
along a shared axis; it is a difference in what kind of thing was written
to. A model of orthographic stabilization that wants to explain why some
historical changes to a script are recoverable in a way others are not needs
a depth axis, not only a cost axis.

\end{document}
